\documentclass{article}
\usepackage[T1]{fontenc}
\usepackage{lmodern}
\usepackage{graphicx}
\usepackage{amsmath,amssymb}
\usepackage[a4paper,margin=2cm]{geometry}
\usepackage[absolute,overlay]{textpos}
\setlength{\TPHorizModule}{1mm}
\setlength{\TPVertModule}{1mm}
\usepackage[indonesian]{babel}
\usepackage{hyperref}
\setlength{\emergencystretch}{2em}
% unit: Halaman 38

\begin{document}

% === Halaman 38 (llm) ===

\begin{itemize}
    \item Kelebihan Vigenere Cipher: huruf plainteks yang sama tidak selalu dienkripsi menjadi huruf ciphteks yang sama pula, bergantung huruf kunci yang digunakan.

    Contoh: pada contoh di atas, huruf plainteks $\mathbb{T}$ dapat dienkripsi menjadi $\mathbb{L}$ atau $\mathbb{H}$, dan huruf cipherteks $\mathbb{V}$ dapat merepresentasikan huruf plainteks $\mathbb{H}$, $\mathbb{I}$, dan $\mathbb{X}$

    \item Hal di atas merupakan karakteristik dari \textit{cipher} abjad-majemuk: setiap huruf cipherteks dapat memiliki kemungkinan banyak huruf plainteks.

    \item Bandingkan dengan \textit{cipher} abjad-tunggal, setiap huruf cipherteks selalu menggantikan huruf plainteks tertentu.
\end{itemize}

\end{document}
